% Labels of main_SR.tex, exported from main_SR.aux by
% export_xr_labels.py. Read by preamble_SR.tex only when the .aux
% itself is absent (Overleaf drops .aux files from uploaded ZIPs).
% Do not edit: regenerate with .\build_SR.ps1.
\newlabel{fig:framework}{{1}{4}{\textbf {Cross-domain medical prediction as an imperfect translation, and the two channels over which it is decomposed.} \textbf {a}, A \emph {source} domain of tissue morphology, whose vocabulary is the six ANORAK growth patterns, is mapped by a multiple-instance translator onto a \emph {target} domain of somatic driver mutations; the per-target \auroc {} measures translational faithfulness averaged over the cohort. \textbf {b}, The two channels the aggregator receives. The morphology channel carries \ctranspath {} tile embeddings $\mathbf {h}_i$ into a gated-attention aggregator (\piabmil {}); the composition channel carries the $6$-d fuzzy pattern probabilities $\mathbf {p}_i$ into a $2$-additive Choquet integral (\fcmil {}). The two are disjoint until the classifier head, which is what makes \cref {eq:pid} well posed as a statement about two sources: neither channel can route information through the other, so an empty pattern-side atom cannot be an artefact of the architecture. Only the vocabulary and the encoder change when the framework moves to glioma MRI or to the electrocardiogram (\cref {fig:pipeline_glioma,fig:pipeline_ecg})}{figure.caption.3}{main_SR.pdf}}
\newlabel{fig:framework@cref}{{[figure][1][]1}{[1][4][]4}{}{}{main_SR.pdf}}
\newlabel{tab:mutation_luad_resultssec}{{1}{6}{\textbf {Mutation prediction on $458$ TCGA-LUAD slides from $444$ patients, five seeds, two outer protocols.} Pooled out-of-fold \auroc {}, mean $\pm $ SD across seeds, under nested cross-validation. B2~$=$~\abmil {} on \ctranspath {} embeddings (the ceiling baseline); B3~$=$~\abmil {} on pattern probabilities only; PI~$=$~\piabmil {} concat fusion; FC~$=$~\fcmil {} $2$-additive Choquet. Best per gene in bold. The B3$-$B2 rows give the hierarchical bootstrap interval over seeds and patients, site-clustered in \textbf {b}; bold marks the genes on which it excludes zero---\textit {TP53}, \textit {EGFR} and \textit {STK11} in \textbf {a}, and only \textit {TP53} and \textit {EGFR} in \textbf {b}, where the \textit {STK11} interval reaches $+0.001$---and these are the only resolved contrasts in the table. No fusion contrast (\piabmil {}$-$B2, \fcmil {}$-$B2, \fcmil {}$-$\piabmil {}) resolves on any gene under either protocol, and no equivalence is established either, at the $\pm 0.02$ margin the regime reads require (\cref {tab:luad_contrasts,tab:design_power})}{table.caption.6}{main_SR.pdf}}
\newlabel{tab:mutation_luad_resultssec@cref}{{[table][1][]1}{[1][5][]6}{}{}{main_SR.pdf}}
\newlabel{tab:luad_contrasts}{{2}{6}{\textbf {Every contrast the regime rule reads, with its interval.} For each gene and outer protocol, the four between-condition contrasts of the noise-aware rule (Methods): \piabmil {}$-$B2 decides additive complementarity (ii-a); \fcmil {}$-$B2 together with \piabmil {}$-$B2 must fall inside $\pm 0.02$ for an equivalence read of regime (i) or (iii); \fcmil {}$-$\piabmil {} decides the non-additive sub-form (ii-b); and B3$-$B2 measures what the vocabulary omits. Each cell gives the pooled out-of-fold point contrast, its $95\%$ hierarchical-bootstrap interval over seeds and patients ($2{,}000$ draws; site-clustered in \textbf {b}), and the number of the five seeds on which the point contrast is positive. Bold marks an interval excluding zero; a superscript $\approx $ would mark an interval inside the $\pm 0.02$ equivalence margin, and no cell earns one. Of $48$ contrasts, five exclude zero and all five are B3$-$B2; none of the $36$ fusion contrasts does, and none is equivalent to zero, which is why every LUAD regime read is reported as directional and unadjudicated. Generated from the same bootstrap output as the B3$-$B2 rows of \cref {tab:mutation_luad_resultssec}}{table.caption.7}{main_SR.pdf}}
\newlabel{tab:luad_contrasts@cref}{{[table][2][]2}{[1][5][]6}{}{}{main_SR.pdf}}
\newlabel{tab:crossdomain_merged}{{3}{10}{\textbf {The procedure run unchanged in two further modalities and across cohorts.} \textbf {a}, Glioma MRI\,$\rightarrow $\,MGMT promoter methylation with the VASARI vocabulary: the verdict is resolved at full power and the regime is read from levels. \textbf {b}, Electrocardiogram\,$\rightarrow $\,physiological state with the graded SCP-ECG statements: the one cohort in which the equivalence condition of the rule is met. \textbf {c}, Cohort transfer from TCGA-LUAD to CPTAC-LUAD. Only the vocabulary and the encoder change between panels; the four conditions, the directional regime rule and the Conclusive/Inconclusive threshold are identical throughout, which is the precise sense in which the procedure rather than the model transfers. The contrasts from which the regimes are read are \cref {tab:crossdomain_contrasts}}{table.caption.14}{main_SR.pdf}}
\newlabel{tab:crossdomain_merged@cref}{{[table][3][]3}{[1][10][]10}{}{}{main_SR.pdf}}
\newlabel{tab:crossdomain_contrasts}{{4}{11}{\textbf {Every contrast the regime rule reads in the cross-domain arms, with its interval.} The counterpart of \cref {tab:luad_contrasts} for the three full-power targets of \cref {tab:crossdomain_merged}: \textbf {a}, MGMT, where no contrast resolves and the regime is read from levels; \textbf {b}, \textbf {c}, sex and age from the electrocardiogram, where both fusion-versus-encoder contrasts fall inside the $\pm 0.02$ equivalence margin and regime~(i) is adjudicated}{table.caption.15}{main_SR.pdf}}
\newlabel{tab:crossdomain_contrasts@cref}{{[table][4][]4}{[1][10][]11}{}{}{main_SR.pdf}}
\newlabel{eq:fuzzy_measure}{{1}{17}{FC-MIL (2-additive fuzzy Choquet)}{equation.1}{main_SR.pdf}}
\newlabel{eq:fuzzy_measure@cref}{{[equation][1][]1}{[1][17][]17}{}{}{main_SR.pdf}}
\newlabel{eq:choquet_integral}{{2}{17}{FC-MIL (2-additive fuzzy Choquet)}{equation.2}{main_SR.pdf}}
\newlabel{eq:choquet_integral@cref}{{[equation][2][]2}{[1][17][]17}{}{}{main_SR.pdf}}
\newlabel{eq:pid}{{3}{18}{Decomposition}{equation.3}{main_SR.pdf}}
\newlabel{eq:pid@cref}{{[equation][3][]3}{[1][18][]18}{}{}{main_SR.pdf}}
\newlabel{eq:modality_ceiling}{{4}{19}{The ceiling is the modality, not the encoder}{equation.4}{main_SR.pdf}}
\newlabel{eq:modality_ceiling@cref}{{[equation][4][]4}{[1][19][]19}{}{}{main_SR.pdf}}
\newlabel{eq:conclusive_verdict}{{5}{27}{Conclusive verdict criterion}{equation.5}{main_SR.pdf}}
\newlabel{eq:conclusive_verdict@cref}{{[equation][5][]5}{[1][27][]27}{}{}{main_SR.pdf}}
